\chapter{СУДАЛГААНЫ ҮНДЭСЛЭЛ}

\section{Асуудлын тодорхойлолт}

Их сургуулийн оюутнууд өдөр бүр олон төрлийн академик мэдээлэлтэй харьцдаг. Хичээлийн хуваарь, шалгалтын хуваарь, дүн, кредит, төлбөрийн мэдээлэл, мэдэгдэл, баримт бичиг зэрэг өгөгдлүүд нь SISI порталын янз бүрийн цэс, дэд хуудсанд тархан байрладаг. Энэ нь хэрэглэгчийг нэг энгийн мэдээлэл авахын тулд олон товшилт хийх, хэд хэдэн дэлгэц дамжих шаардлагатай болгодог.

Жишээлбэл, оюутан ``маргаашийн хичээл хэдэн цагт эхлэх вэ'' гэсэн энгийн асуултад хариулахын тулд эхлээд порталд нэвтэрч, хуваарийн цэс рүү орж, тухайн өдрийн хичээлүүдийг олж, цагийг шалгах шаардлагатай. Үүнтэй адил ``миний GPA хэд вэ'', ``төлбөрийн үлдэгдэл хэд вэ'', ``дараагийн шалгалт хэзээ вэ'' зэрэг түгээмэл асуултууд бүгд олон алхам шаарддаг.

Энэ асуудал нь зөвхөн интерфейсийн төвөгтэй байдлаас үүдээгүй. Учир нь порталын архитектур нь өгөгдлийг функцээр нь салгасан бөгөөд хэрэглэгчийн ``зорилго''-ыг ойлгож өгдөг ухаалаг давхарга байхгүй. Одоогийн портал нь ``хуваарь харуулах'', ``дүн харуулах'' гэсэн тусдаа функцуудтай боловч ``миний энэ долоо хоногт хамгийн ачаалалтай өдөр аль вэ'' эсвэл ``миний сайжруулах шаардлагатай хичээлүүд юу вэ'' гэсэн илүү өндөр түвшний асуултад хариулах чадваргүй.

\section{Автоматжуулалтын хэрэгцээ}

Оюутнууд зөвхөн мэдээлэл уншаад зогсохгүй, уг мэдээллийг өөр системүүд рүү дамжуулах, боловсруулах хэрэгцээтэй байдаг. Тухайлбал, хичээлийн хуваариа Google Calendar-д оруулах, даалгаврын deadline-ийг todo app-д нэмэх, дүнгүүдийг Excel файл болгож хадгалах зэрэг ажлууд өдөр тутам давтагддаг. Эдгээр нь нэг бүрчлэн гүйцэтгэхэд хялбар боловч нийлээд ихээхэн цаг хугацаа, анхаарал шаарддаг.

Багш нарын хувьд ч мөн адил асуудал байдаг. Ялангуяа семестр бүрийн эцэст олон оюутны дүнг Excel файлаас системд гараар оруулах нь маш их цаг хугацаа шаардсан, алдаа гарах магадлал өндөр ажил юм. Нэг хичээлд 50-100 оюутан байх тохиолдолд энэ үйлдлийг автоматжуулах нь ихээхэн ач холбогдолтой.

\section{Хиймэл оюуны одоогийн чадамж ба хязгаарлалт}

Сүүлийн жилүүдэд хэлний загварууд (Large Language Models) гайхалтай хурдацтай хөгжиж, текст ойлгох, боловсруулах, үүсгэх чадварын хувьд хүний түвшинд хүрсэн байна. GPT-4, Claude, Gemini зэрэг загварууд нь нарийн төвөгтэй асуултад хариулах, код бичих, баримт бичиг боловсруулах, шинжилгээ хийх зэрэг даалгаврыг өндөр чанартай гүйцэтгэж чаддаг.

Гэсэн хэдий ч эдгээр загварууд нь ``хаалттай орчин''-д ажилладаг. Тэдгээр нь сургалтын үеийн мэдлэг дээр тулгуурлан хариулт үүсгэдэг бөгөөд бодит цагийн өгөгдөлд хандах, гаднах системтэй харилцах, хэрэглэгчийн өмнөөс үйлдэл гүйцэтгэх чадваргүй. Өөрөөр хэлбэл, загварын ``оюун ухаан'' хангалттай боловч бодит ертөнцтэй харилцах ``гар, хөл'' нь дутагдалтай.

Энэ хязгаарлалтыг шийдэх хамгийн үр дүнтэй арга нь загварын оюун ухааныг нэмэгдүүлэх бус, түүнд зохих хэрэгслүүдийг олгох явдал юм. Загварын параметрийн тоог нэмэгдүүлэх, илүү их өгөгдөл дээр сургах зэрэг incremental сайжруулалтаас илүүтэйгээр, загварт бодит системүүдтэй харилцах интерфейс хангах нь шинэ түвшний чадамж нээж өгдөг.

\section{Хувийн сэдэлт}

Энэхүү төслийг хэрэгжүүлэх хувийн сэдэлт нь дээрх онолын үндэслэлээс гадна практик туршлагаас үүдэлтэй. Би өөрөө NUM-ийн оюутан болохын хувьд SISI порталыг өдөр тутам ашигладаг бөгөөд түүний давуу болон сул талуудыг сайн мэддэг. Энэ нь надад бодит хэрэгцээнд тулгуурласан систем загварчлах, хөгжүүлэх боломжийг олгож байна.

Үүнээс гадна хиймэл оюуны загварт хэрэгсэл олгох (tool-use) чиглэлийг гүнзгий судлах, туршиж үзэх хүсэлтэй байсан. Загварын суурь оюун ухаан нь аль хэдийн хангалттай түвшинд хүрсэн бөгөөд түүнд зөв хэрэгслүүдийг олгосноор бодит амьдрал дээр нөлөө үзүүлэх систем бүтээх боломжтой гэдэгт итгэдэг. Энэхүү итгэл үнэмшлийг өөрт ойр, сайн мэддэг системээр туршиж баталгаажуулах нь энэ төслийн нэг зорилго юм.

\section{Асуудлын томьёолол}

Дээрх үндэслэлд тулгуурлан энэхүү судалгааны асуудлыг дараах байдлаар томьёолж байна: SISI портал болон холбогдох бусад системүүдийн өгөгдөл, үйлчилгээг найдвартай, аюулгүй, өргөтгөх боломжтой хэлбэрээр байгалийн хэлний интерфейстэй хэрхэн холбож, хэрэглэгчийн хүсэлтийг бодит системийн үйлдэл болгон хөрвүүлэх вэ?

Энэ асуудлыг шийдэхэд зөвхөн хэлний интерфейсийн өнгөц давхарга хангалтгүй. Систем нь бодит цагийн өгөгдөл ашиглах, хэрэглэгчийг дахин нэвтрүүлэхгүйгээр одоогийн сессийг аюулгүй ашиглах, унших үйлдэл болон бичих үйлдлийг ялгаж удирдах, олон системийн хооронд уялдаа хангах, мөн оюутны болон багшийн өөр өөр хэрэглээний хувилбарт тохирох модульчлагдсан бүтэцтэй байх шаардлагатай.

Иймээс энэхүү дипломын ажлын судалгааны гол чиглэл нь агентын онол, Model Context Protocol стандартыг ашигласан архитектур, мөн бодит интеграцын түвшний системийн загварчлал болно.
